\scanpage{95}
% archive foliation: 94
Donc pour la suite on se donne d'abord l'extension
\[
0\longrightarrow\Omega\longrightarrow\mathcal{L}\longrightarrow\underline{\mathcal{O}}_X\longrightarrow0
\]
\margin{On suppose $2$ inv. sur $X$}
($\Omega$ inversible), puis l'op. diff. d'ordre $\leq1$
\[
\Omega\xrightarrow{c_\Omega}\mathcal{L}\otimes\underline{\omega}\qquad\underline{\omega}=\Omega^1_{X/S}
\]
tel que

a) le composé avec $\mathcal{L}\otimes\underline{\omega}\longrightarrow\underline{\omega}$ soit \deleted{\unsure{un}} $\Omega\longrightarrow\underline{\omega}$ linéaire et un \uline{iso} (qu'on notera $-2i$, $i:\Omega\simeq\underline{\omega}$ qu'on utilisera parfois pour identifier)
\margin{\unsure{cela a l'air vraiment vicieux vrai}}

b) \deleted{\unsure{L'hom.}} On a (si $w\in\Gamma\Omega$, $f\in\Gamma\underline{\mathcal{O}}_X$)
\[
c_\Omega(fw)-f\,c_\Omega(w)=w\otimes df
\]
($w$ fixé, \unsure{c'est bien} \uline{lin.} $\Omega\longrightarrow\Omega\otimes\omega$)

[Ces conditions définissent le \unsure{système} de l'ext. $\mathcal{L}$ et de l'op. diff $c_\Omega$ à isom unique près -- et on peut tjrs en trouver, si $\underline{\omega}=\Omega^1_{X/S}$ est bien inversible)

On forme $\mathcal{H}$, et on considère les connexions \inserted{\unsure{\(\mathcal{H}\)}} d'alg. de Lie \deleted{\unsure{\(\mathcal{H}\)}} sur $\mathcal{H}$ qui prolongent $c_\Omega$. Elles forment les sections d'un torseur sous $\unsure{\Omega^2}\otimes\underline{\omega}\simeq\underline{\omega}^{\otimes\unsure{2|3}}$. \deleted{On se demande quelles}

\margin{La donnée de $c_\Omega$ revient à celle d'un iso. des voisinages infinitésimaux d'ordre $2$ de $\sigma(X)$ dans $P$ et de $\Delta_X(X)$ dans $X\times_S X$. $\unsure{\longrightarrow}$ celle de $c_{\mathcal{H}}$ les \deleted{\unsure{prolongeant à un iso des}} voisinages d'ordre $3$ les prolongeant (et alors, un iso des \unsure{vois. form.})}

Faire le lien avec la situation $(P,\sigma,c)$ et les données formelles
\[
\underline{\mathcal{O}}_{\widehat P}^{\mathrm{con}}\simeq P^\infty_{X/S}
\]
(si $X/S$ de \unsure{f.t.} lisse --).
